% ECONOMICS_PROBLEM: {"id": "D73-385", "title": "Durable anticorruption policy", "jel_code": "D73", "source_id": "OP-0383"}
% !TeX program = xelatex
\documentclass[11pt,a4paper]{article}
\usepackage[margin=23mm]{geometry}
\usepackage{fontspec}
\setmainfont{Latin Modern Roman}
\usepackage{amsmath,amssymb,mathtools}
\usepackage{xcolor,enumitem,booktabs,longtable,array,xurl,hyperref}
\definecolor{muted}{HTML}{53636C}
\hypersetup{colorlinks=true,urlcolor=blue,linkcolor=blue}
\setlength{\parindent}{0pt}
\setlength{\parskip}{5pt}
\newcommand{\R}{\mathbb R}
\newcommand{\E}{\mathbb E}
\newcommand{\N}{\mathbb N}
\newcommand{\ind}{\mathbf 1}
\newcommand{\Var}{\operatorname{Var}}
\newcommand{\Cov}{\operatorname{Cov}}
\newcommand{\supp}{\operatorname{supp}}
\DeclareMathOperator*{\argmin}{arg\,min}
\DeclareMathOperator*{\argmax}{arg\,max}
\newcommand{\entrymeta}[3]{{\sffamily\small\color{muted}\textbf{\texttt{#1}}\quad|\quad #2\par #3\par}}
\newcommand{\Problemref}[1]{\texttt{#1}}
\newcommand{\JELref}[2]{\texttt{#2} (JEL \texttt{#1})}
\begin{document}
\section*{Durable anticorruption policy}
\textbf{Problem ID:} \texttt{D73-385}\quad \textbf{JEL:} \texttt{D73}\par
% BEGIN ECONOMICS STATEMENT
\label{problem:OP-0383}
{\sffamily\small\color{muted}Original subject group: Political economy and institutions\par}
\label{definitions:problem:30007665}
\textbf{Audit classification:} Published research agenda. Primary manuscript and metadata checked. Learning, substitution and long-run evaluation are explicit gaps; harm weights and measurement identification are editorial.
\subsection*{Definitions and precise research target}
\textbf{Complete editorial problem.} Fix a public organization and a policy bundle \(p\) comprising audits, sanctions, transparency, competition, and personnel incentives. Let \(C^k_{jt}(p)\) be verified corruption through avenue \(k\) in unit \(j\) at date \(t\), and let \(\lambda_k\ge0\) be declared weights capturing losses rather than simply counts of detected offenses. Define total corruption harm \(C_{jt}(p)=\sum_k\lambda_kC^k_{jt}(p)\). The main estimand is \(\tau_C(T,p)=E[C_{jT}(p)-C_{jT}(0)]\), for a horizon permitting officials to learn and substitute across avenues. Estimate service output, enforcement cost, selection into office, and citizen welfare alongside \(\tau_C\). Audits may change detection probabilities, so measured offenses require an explicit measurement model or independent validation. A policy that lowers the targeted avenue while increasing another need not reduce total harm. Use credible bundle and follow-up variation to distinguish deterrence, adaptation, and personnel selection; describe interference across agencies. Determine which combinations remain effective after officials understand the rules. The source explicitly requests short- and long-run evidence under strategic adaptation. The avenue weights, outcome definition, and comparative policy objective are our formalization of that research agenda.

\textbf{Definitions and admissible scope.} Choose mutually exclusive corruption avenues or explicitly remove double counting. \(C^k_{jt}\) is the true amount or incidence of avenue \(k\), while detected cases \(D^k_{jt}\) are generated by sensitivity, false-positive rates, and audit intensity that can depend on policy. Harm weights \(\lambda_k\) have a fixed welfare or currency interpretation; detection counts are not harm. Long-run effectiveness means \(E[C_{jT}(p)-C_{jT}(0)]<0\) at a declared horizon after learning, accompanied by service outcomes and cost. If every \(\lambda_k\ge0\), reducing one avenue is insufficient when others increase. The assignment law specifies interagency interference and officials' opportunities to learn and move between posts. Let the corruption avenue set be finite, with \(\lambda_k\ge0\) and \(\sum_k\lambda_k>0\), finite expected harms, a fixed organization-unit law, and finite horizon \(T\). The complete corruption definition includes displacement to uncovered units or states whether it lies outside the target. Fix independent validation or explicit sensitivity/false-positive bounds before comparing detected offenses. The target is the joint identified set of total-harm, service, and cost responses under a declared finite policy menu and strategic response/assignment law. A cost-effective ordering requires a separately stated welfare criterion; a negative \(\tau_C\) alone need not maximize citizen welfare.
\subsection*{Variants and assumption changes}
\begin{enumerate}
\item Detection-adjusted deterrence (editorial): hold independent audit sensitivity fixed while varying sanctions or regular-audit intensity. Identify harm changes separately from discovery changes.
\item Adaptation and selection (source-motivated): collect repeated measures across avenues and personnel cohorts after officials learn the policy. Compare changing audit patterns and recruitment rules to identify durable substitution and entry effects.
\end{enumerate}
\subsection*{References and source audit}
\textbf{Support and access.} Primary manuscript and metadata checked. Learning, substitution and long-run evaluation are explicit gaps; harm weights and measurement identification are editorial. \textbf{Locator:} February 24, 2012 manuscript, Section 4 pp. 35--37 and conclusion pp. 38--39.
\begin{enumerate}\raggedright
\item\hypertarget{definitions:source:30007665:1}{} Benjamin A. Olken; Rohini Pande. 2012. \emph{Corruption in Developing Countries}. February 24, 2012 manuscript, Section 4, printed pp. 35--37; conclusion, pp. 38--39.
\url{https://economics.mit.edu/research/publications/corruption-developing-countries}.
DOI: \url{https://doi.org/10.1146/annurev-economics-080511-110917}.
\end{enumerate}

\subsection*{Common conventions: Source mathematical conventions}

These conventions apply to each entry unless explicitly replaced.

\textbf{Models and identification.} A primitive model \(m\) specifies all probability laws, feasible choices, information, equilibrium and measurement rules used by its target. The admissible class \(\mathcal M\) is fixed independently of outcomes. If \(P_O(m)\) is its observable probability law and \(\tau(m)\) the target, its identified set at observable law \(P\) is
\[
 \mathcal I_\tau(P)=\{\tau(m):m\in\mathcal M,\ P_O(m)=P\}.
\]
Point identification means that this set is a singleton. An empty set signals incompatibility of data and assumptions. Coordinate bounds need not describe the joint set. Bounds are sharp only if every retained target is attainable or arbitrarily approachable by compatible models. Finite bounds require explicit restrictions on unobserved quantities; unrestricted confounding, hidden wealth, extrapolation or arbitrary equilibrium selection can yield uninformative sets.

\textbf{Causal interventions.} A potential outcome \(Y_i(p)\) is the outcome under a complete policy regime \(p\), including timing, financing, information and market spillovers. The target population and units are fixed before treatment. With interference, use the complete assignment vector or a justified exposure mapping; individual no-interference notation alone cannot identify an economy-wide intervention. A difference of mean potential outcomes does not require the cross-world joint distribution. Conditioning on post-treatment migration, survival, enrollment or employment changes the estimand and needs separate justification.

\textbf{Statistical statements.} An estimator, confidence set, forecast or test specifies a sampling law, model class and loss/coverage criterion. Uniform \((1-\alpha)\)-coverage means \(\inf_{m\in\mathcal M}\Pr_m\{\tau(m)\in C_n\}\geq1-\alpha\), or an explicitly asymptotic version. The whole parameter space trivially has coverage, so informativeness requires a stated width, risk or power objective. A forecasting improvement is relative to a named benchmark, information set and evaluation population.

\textbf{Equilibrium and optimization.} An equilibrium comprises feasible plans, optimality, beliefs where needed, budget constraints and all market/resource clearing. The primitives or a stated benchmark define each correspondence \(\mathcal E(p;m)\). Nonemptiness is assumed or proved, never inferred just from compact policy sets. A selected-equilibrium target uses a declared selection rule; a robust target quantifies over all permitted equilibria. Use suprema or infima unless attainment is established. An \(\varepsilon\)-optimal policy is feasible and within \(\varepsilon\) of the finite optimum value. Report an empty feasible set separately.

\textbf{Welfare and accounting.} Normative utility, interpersonal normalization, social weights, numeraire, discounting and the use of public receipts are primitives. Behavioral utility need not coincide with the welfare criterion. Household welfare includes ownership income and rents once; adding profits again to compensating variation generally double-counts them. A monetary equivalent is defined by an explicit transfer and price convention, with monotonicity and range conditions. Average effects, mechanism contributions, transfers and welfare losses are different targets. Infinite sums require convergence and dynamic feasible plans require terminal or no-Ponzi conditions.

\textbf{Source versions.} A working-paper date, revision date and journal publication year are distinguished. A DOI for a journal version is not automatically the DOI of an earlier manuscript. Locators refer to the stated version and printed pages where specified. A metadata-only check does not verify a claim about a full-text conclusion. Every entry's source subsection narrows the provenance claim accordingly.
% END ECONOMICS STATEMENT
\end{document}
